% Répartition du budget de latence par étape (figure 4.2)
\begin{tikzpicture}[x=1cm, y=1cm]

  % Échelle : 1 cm pour 100 ms
  \def\S{0.01}

  \tikzset{
    seg/.style={draw=black},
    lg/.style={font=\sffamily\tiny, anchor=west}
  }

  % ── Axe des abscisses ─────────────────────────────────────
  \draw[->] (0,-2.9) -- (21.5,-2.9);
  \foreach \ms in {0,250,500,750,1000,1250,1500,1750,2000}{
    \draw ({\ms*\S},-2.9) -- ({\ms*\S},-3.05);
    \node[font=\sffamily\tiny, anchor=north] at ({\ms*\S},-3.08) {\ms};
  }
  \node[font=\sffamily\tiny\itshape, anchor=north] at (10.75,-3.60) {millisecondes};

  % ── Barre 1 : échange sans recherche documentaire ─────────
  % motifs distincts par étape, sans couleur
  \def\ya{-0.35}
  \def\hb{0.72}
  \fill[figGrisClair, seg]   (0,\ya) rectangle ({250*\S},\ya-\hb);
  \fill[white, seg]          ({250*\S},\ya) rectangle ({450*\S},\ya-\hb);
  \fill[figGrisTresClair, seg]({450*\S},\ya) rectangle ({495*\S},\ya-\hb);
  \fill[figGrisMoyen, seg]   ({495*\S},\ya) rectangle ({1045*\S},\ya-\hb);
  \fill[white, seg]          ({1045*\S},\ya) rectangle ({1245*\S},\ya-\hb);
  \fill[figGrisClair, seg]   ({1245*\S},\ya) rectangle ({1320*\S},\ya-\hb);
  \node[font=\sffamily\tiny, anchor=east] at (-0.15,\ya-0.36) {Sans recherche};
  \node[font=\sffamily\tiny\bfseries, anchor=west] at ({1320*\S+0.15},\ya-0.36) {1,0 à 1,6 s};

  % ── Barre 2 : échange avec recherche documentaire ─────────
  \def\yb{-1.45}
  \fill[figGrisClair, seg]   (0,\yb) rectangle ({250*\S},\yb-\hb);
  \fill[white, seg]          ({250*\S},\yb) rectangle ({450*\S},\yb-\hb);
  \fill[figGrisTresClair, seg]({450*\S},\yb) rectangle ({495*\S},\yb-\hb);
  \fill[pattern=north east lines, seg] ({495*\S},\yb) rectangle ({770*\S},\yb-\hb);
  \fill[figGrisMoyen, seg]   ({770*\S},\yb) rectangle ({1320*\S},\yb-\hb);
  \fill[white, seg]          ({1320*\S},\yb) rectangle ({1520*\S},\yb-\hb);
  \fill[figGrisClair, seg]   ({1520*\S},\yb) rectangle ({1595*\S},\yb-\hb);
  \node[font=\sffamily\tiny, anchor=east] at (-0.15,\yb-0.36) {Avec recherche};
  \node[font=\sffamily\tiny\bfseries, anchor=west] at ({1595*\S+0.15},\yb-0.36) {1,2 à 2,0 s};

  % ── Seuil de perception ───────────────────────────────────
  \draw[dashed, line width=0.7pt] ({1500*\S},0.35) -- ({1500*\S},-2.9);
  \node[font=\sffamily\tiny\itshape, anchor=south, align=center] at ({1500*\S},0.42)
       {seuil au delà duquel\\l'attente devient perceptible};

  % ── Légende ───────────────────────────────────────────────
  \def\yl{-4.35}
  \foreach \x/\y/\st/\lbl in {
      0/0/{figGrisClair}/{Détection de fin d'énoncé et gigue réseau},
      10.6/0/{white}/{Transcription et synthèse vocale},
      0/1/{figGrisTresClair}/{Assemblage du contexte},
      10.6/1/{figGrisMoyen}/{Raisonnement}}
    {
      \fill[\st, draw=black] (\x,{\yl-\y*0.58}) rectangle (\x+0.42,{\yl-\y*0.58+0.30});
      \node[lg] at (\x+0.52,{\yl-\y*0.58+0.15}) {\lbl};
    }
  \fill[pattern=north east lines, draw=black] (0,\yl-1.16) rectangle (0.42,\yl-0.86);
  \node[lg] at (0.52,\yl-1.01)
       {Récupération documentaire, présente uniquement lorsque la demande l'appelle};

\end{tikzpicture}
